% Continuation fragment for GAGA source nos. 9--12. Insert after the
% analytification fragment and before the proof fragment for nos. 13--17.

\section{Analytification of modules}
\label{section-analytification-modules}

\begin{definition}
\label{definition-analytification-module}
\begin{reference}
\cite[\S 3, no. 9, Definition 2]{GAGA}
\end{reference}
Let $X$ be an algebraic variety over $\mathbf C$ and let
$a_X:X^{\mathrm{an}}\to X$ be the canonical morphism of locally ringed
spaces of Proposition \ref{proposition-analytification}. For an
$\mathcal O_X$-module $\mathcal F$, its {it analytification} is
$$
\mathcal F^{\mathrm{an}}=a_X^*\mathcal F
=\mathcal O_{X^{\mathrm{an}}}
\otimes_{a_X^{-1}\mathcal O_X}a_X^{-1}\mathcal F.
$$
There is a canonical map of $a_X^{-1}\mathcal O_X$-modules
$$
\eta_{\mathcal F}:a_X^{-1}\mathcal F\longrightarrow
\mathcal F^{\mathrm{an}},\qquad s\longmapsto1\otimes s.
$$
Every $\mathcal O_X$-linear map $\mathcal F\to\mathcal G$ induces an
$\mathcal O_{X^{\mathrm{an}}}$-linear map
$\mathcal F^{\mathrm{an}}\to\mathcal G^{\mathrm{an}}$, functorially, and
$\mathcal O_X^{\mathrm{an}}=\mathcal O_{X^{\mathrm{an}}}$ canonically.
\end{definition}

\begin{lemma}
\label{lemma-analytification-module-stalks}
Let $x\in X(\mathbf C)=X^{\mathrm{an}}$. There are canonical identifications
$$
(\mathcal F^{\mathrm{an}})_x
=\mathcal F_x\otimes_{\mathcal O_{X,x}}
\mathcal O_{X^{\mathrm{an}},x}
$$
under which $(\eta_{\mathcal F})_x$ is the map $m\mapsto1\otimes m$.
\end{lemma}

\begin{proof}
Stalks commute with inverse image and tensor product. Since $x$ is a closed
point, $(a_X^{-1}\mathcal F)_x=\mathcal F_x$. The displayed identification
and the description of $\eta_{\mathcal F}$ follow directly from Definition
\ref{definition-analytification-module}.
\end{proof}

\begin{lemma}
\label{lemma-analytification-exact-faithful-coherent}
\begin{reference}
\cite[\S 3, no. 9, Proposition 10]{GAGA}
\end{reference}
Let $X$ be an algebraic variety over $\mathbf C$.
\begin{enumerate}
\item The functor $\mathcal F\mapsto\mathcal F^{\mathrm{an}}$ on
$\mathcal O_X$-modules is exact.
\item The map $\eta_{\mathcal F}:a_X^{-1}\mathcal F\to
\mathcal F^{\mathrm{an}}$ is injective for every $\mathcal F$.
\item If $\mathcal F$ is coherent, then $\mathcal F^{\mathrm{an}}$ is a
coherent analytic module.
\end{enumerate}
\end{lemma}

\begin{proof}
The stalk formula of Lemma \ref{lemma-analytification-module-stalks} and
faithful flatness in Theorem
\ref{theorem-analytification-local-faithfully-flat} prove (1); compare
Modules, Lemma \ref{modules-lemma-pullback-flat}. The same stalk formula and
Algebra, Lemma
\ref{algebra-lemma-faithfully-flat-universally-injective} prove (2).

For (3), a coherent module on the locally Noetherian scheme $X$ is locally
finitely presented. Pullback preserves finite presentations. The analytic
structure sheaf is coherent by Theorem
\ref{theorem-oka-cartan-coherence}; hence a finitely presented analytic
module is coherent by Modules, Lemma
\ref{modules-lemma-coherent-structure-sheaf}.
\end{proof}

\begin{remark}[Analytification of an ideal]
\label{remark-analytification-ideal-sheaf}
If $\mathcal J\subset\mathcal O_X$ is a quasi-coherent ideal, exactness
identifies $\mathcal J^{\mathrm{an}}$ with the ideal of
$\mathcal O_{X^{\mathrm{an}}}$ generated by the image of
$a_X^{-1}\mathcal J$. For the ideal of a reduced closed subvariety this is
also Proposition \ref{proposition-analytification-closed-ideal}.
\end{remark}


\section{Closed immersions and analytification}
\label{section-analytification-closed-immersions}

\begin{lemma}
\label{lemma-analytification-pushforward-closed-immersion}
\begin{reference}
\cite[\S 3, no. 10, Proposition 11]{GAGA}
\end{reference}
Let $i:Z\to X$ be a closed immersion of algebraic varieties over
$\mathbf C$, and let $\mathcal F$ be a coherent $\mathcal O_Z$-module. There
is a canonical isomorphism
$$
(i_*\mathcal F)^{\mathrm{an}}
\longrightarrow i^{\mathrm{an}}_*(\mathcal F^{\mathrm{an}}).
$$
\end{lemma}

\begin{proof}
The commutative analytification square and adjunction give the canonical
map. Both sides vanish away from $Z^{\mathrm{an}}$. At
$z\in Z(\mathbf C)$ the source stalk is
$$
\mathcal F_z\otimes_{\mathcal O_{X,z}}
\mathcal O_{X^{\mathrm{an}},z},
$$
whereas the target stalk is
$$
\mathcal F_z\otimes_{\mathcal O_{Z,z}}
\mathcal O_{Z^{\mathrm{an}},z}.
$$
Proposition \ref{proposition-analytification-closed-ideal} gives
$$
\mathcal O_{Z^{\mathrm{an}},z}
=\mathcal O_{Z,z}\otimes_{\mathcal O_{X,z}}
\mathcal O_{X^{\mathrm{an}},z}.
$$
Associativity of tensor product identifies the two stalks, and the
identification is the stalk of the canonical map.
\end{proof}


\section{The cohomology comparison map}
\label{section-analytification-cohomology-map}

\begin{definition}[Cohomology comparison map]
\label{construction-analytification-cohomology-map}
\begin{reference}
\cite[\S 3, no. 11]{GAGA}
\end{reference}
Let $\mathcal F$ be an $\mathcal O_X$-module. Exactness of $a_X^*$ and the
derived adjunction unit give
$$
\mathcal F\longrightarrow Ra_{X,*}a_X^*\mathcal F.
$$
After applying derived global sections and using
$$
R\Gamma(X,Ra_{X,*}\mathcal G)
=R\Gamma(X^{\mathrm{an}},\mathcal G),
$$
we obtain canonical maps
$$
\epsilon^q_{\mathcal F}:
H^q(X,\mathcal F)\longrightarrow
H^q(X^{\mathrm{an}},\mathcal F^{\mathrm{an}}),
\qquad q\geq0.
$$
For $q=0$, this sends a section $s$ to $1\otimes s$. Thus it agrees with
Serre's map obtained from Zariski covers and their analytifications.
\end{definition}

\begin{lemma}
\label{lemma-analytification-cohomology-map-functorial}
The maps $\epsilon^q_{\mathcal F}$ are natural in $\mathcal F$. For every
short exact sequence of $\mathcal O_X$-modules, they form a morphism from
the algebraic long exact cohomology sequence to the analytic long exact
cohomology sequence. In particular, they commute with connecting maps.
\end{lemma}

\begin{proof}
The derived adjunction unit is a natural transformation. Because
analytification is exact by Lemma
\ref{lemma-analytification-exact-faithful-coherent}, it carries short exact
sequences to short exact sequences. Naturality of the associated
distinguished triangles gives the asserted morphism of long exact
sequences.
\end{proof}


\section{The GAGA comparison theorems}
\label{section-gaga-statements}

\begin{theorem}[Cohomology comparison]
\label{theorem-gaga-cohomology}
\begin{reference}
\cite[\S 3, no. 12, Theorem 1]{GAGA}
\end{reference}
Let $X$ be a projective algebraic variety over $\mathbf C$ and let
$\mathcal F$ be a coherent $\mathcal O_X$-module. For every $q\geq0$, the
map
$$
\epsilon^q_{\mathcal F}:H^q(X,\mathcal F)
\longrightarrow H^q(X^{\mathrm{an}},\mathcal F^{\mathrm{an}})
$$
is an isomorphism. In particular,
$$
\Gamma(X,\mathcal F)\longrightarrow
\Gamma(X^{\mathrm{an}},\mathcal F^{\mathrm{an}})
$$
is an isomorphism.
\end{theorem}

\begin{theorem}[Full faithfulness]
\label{theorem-gaga-fully-faithful}
\begin{reference}
\cite[\S 3, no. 12, Theorem 2]{GAGA}
\end{reference}
Let $X$ be a projective algebraic variety over $\mathbf C$ and let
$\mathcal F,\mathcal G$ be coherent $\mathcal O_X$-modules. Analytification
induces a bijection
$$
\mathop{\rm Hom}\nolimits_{\mathcal O_X}(\mathcal F,\mathcal G)
\longrightarrow
\mathop{\rm Hom}\nolimits_{\mathcal O_{X^{\mathrm{an}}}}
(\mathcal F^{\mathrm{an}},\mathcal G^{\mathrm{an}}).
$$
\end{theorem}

\begin{theorem}[Essential surjectivity]
\label{theorem-gaga-essential-surjectivity}
\begin{reference}
\cite[\S 3, no. 12, Theorem 3]{GAGA}
\end{reference}
Let $X$ be a projective algebraic variety over $\mathbf C$. Every coherent
$\mathcal O_{X^{\mathrm{an}}}$-module $\mathcal M$ is isomorphic to
$\mathcal F^{\mathrm{an}}$ for a coherent $\mathcal O_X$-module
$\mathcal F$. The module $\mathcal F$ is unique up to isomorphism; moreover,
every analytic isomorphism between two such analytifications is induced by
a unique algebraic isomorphism.
\end{theorem}

The proofs of these three theorems occupy the following sections.

\begin{remark}[Projectivity is essential]
\label{remark-gaga-projectivity-essential}
For $X=\mathbf A^1_{\mathbf C}$ and $\mathcal F=\mathcal O_X$, the
degree-zero comparison map is
$$
\mathbf C[z]\longrightarrow
\Gamma(\mathbf C,\mathcal O_{\mathbf C}),
$$
whose target is the ring of entire functions. It is not surjective. Thus the
projectivity hypothesis in the comparison theorems cannot simply be
discarded.
\end{remark}

\begin{remark}[The intermediate inverse-image sheaf]
\label{remark-gaga-comparison-factorization}
The comparison map factors canonically as
$$
H^q(X,\mathcal F)
\longrightarrow H^q(X^{\mathrm{an}},a_X^{-1}\mathcal F)
\longrightarrow H^q(X^{\mathrm{an}},\mathcal F^{\mathrm{an}}).
$$
The first arrow need not be an isomorphism. For example, let $X$ be an
irreducible projective variety whose analytification has a nonzero Betti
number $b_q$ with $q>0$, and put $K=\mathbf C(X)$. The sheaf $K$ is flasque
on the irreducible Zariski space $X$, so $H^q(X,K)=0$. Its inverse image is
the constant sheaf with value $K$ on $X^{\mathrm{an}}$, whose degree-$q$
cohomology is a $K$-vector space of dimension $b_q$. This is the obstruction
recorded after the three theorems in \cite[\S 3, no. 12]{GAGA}.
\end{remark}
